\chapter{질량을 가진 약한 상호작용의 교환, 가상 입자, 대칭 깨짐}

표준모형의 군은 페르미온 표에서 허용되는 변환들의 집합으로 보면 훨씬 덜 신비롭게 다가온다. 색 변환은 빨강, 초록, 파랑 사이에서 세로 방향으로 움직인다. 전하를 띤 약한 변환은 쿼크와 렙톤 이중항 안에서 가로 방향으로 움직인다. 전자기 위상 변환은 세 번째 연산을 제공한다. 이에 대응하는 게이지 보손은 대칭을 찾은 뒤에 덧붙인 장식이 아니다. 게이지 보손을 방출하면 대칭의 생성자가 기술하는 것과 같은 상태 변화가 물리적으로 일어난다.

이 관찰을 바탕으로 이번 강의의 핵심 질문을 제기할 수 있다. 약한 상호작용의 꼭짓점 결합은 전자기 상호작용의 꼭짓점 결합과 비슷한 크기인데도, 베타 붕괴는 전형적인 강한 상호작용 과정이나 전자기 과정에 비해 엄청나게 느리다. 여기에 빠진 요소는 전파다. 운동량이 낮을 때 질량이 있는 \(W\)를 교환하면 진폭에 대략 \(1/m_W^2\)의 인자가 곱해진다. 먼저 거리 공간에서, 다음으로 운동량 공간에서 이 설명을 전개하고, 실재 입자와 가상 입자에 관한 긴 논의를 거쳐 마지막으로 자발적 대칭 깨짐의 첫 정성적 모형을 살펴보겠다.

\section{페르미온 표에서 게이지 대칭 읽기}

\subsection{색은 세로 방향으로 작용한다}

여섯 가지 쿼크 맛을 세 약한 상으로 배열하고,
\begin{equation}
  (u,d),\qquad(c,s),\qquad(t,b),
\end{equation}
각 쿼크에 빨강, 초록 또는 파랑의 색 표지를 붙이자. 이렇게 얻은 쿼크 표에는 색 행이 세 개, 맛 열이 여섯 개 있다. 색 \(SU(3)\)은 어느 한 맛 열 안의 세 항목을 서로 섞는다. 예를 들어 한 생성자는 빨간 업 쿼크를 초록 업 쿼크로 바꿀 수 있다. 입자 언어로 같은 작용을 말하면, 보상하는 색과 반색을 지닌 글루온의 방출이 동반된다.

같은 색 변환은 모든 맛에 작용해야 한다. 빨간 업 쿼크를 초록 업 쿼크로 바꾸는 글루온은 빨간 톱 쿼크를 초록 톱 쿼크로 바꾸는 글루온과 같은 종이다. 진정한 대칭 작용은 1세대에만 적용되고 나머지 세대는 무시하지 않는다. 따라서 색 변환은 세로 방향으로 일어나며 보편적으로 작용한다.
\begin{equation}
  q_r\longleftrightarrow q_g\longleftrightarrow q_b,
  \qquad q\in\{u,d,c,s,t,b\}.
\end{equation}
색은 바꾸지만 맛은 바꾸지 않는다.

이것은 게이지 군과 그 게이지 보손 사이의 전형적인 관계다. 무한소 생성자는 상태들 사이의 변환을 기술하며, 이에 대응하는 보손을 방출하거나 흡수하면 그 변환이 물리적 과정으로 구현된다.

\subsection{약한 변환은 가로 방향으로 작용한다}

전하를 띤 약한 상호작용은 다른 방향으로 작용한다. 업형 쿼크는 색을 바꾸지 않은 채 전하를 띤 약한 보손을 방출하고 다운형 짝으로 바뀔 수 있다.
\begin{equation}
  u_b\longrightarrow d_b+W^+.
\end{equation}
정확히 같은 \(W^+\)가
\begin{equation}
  c_g\longrightarrow s_g+W^+,
  \qquad
  t_r\longrightarrow b_r+W^+.
\end{equation}
에서도 일어난다. \(r,g,b\) 표지는 이 과정에서 바뀌지 않는다. 글루온은 표를 가로질러 움직이지 않으며, 전하를 띤 \(W\) 보손은 색 표에서 세로 방향으로 움직이지 않는다.

전자기 위상 대칭 \(U(1)\)을 더하면 모든 전하를 띤 장에 위상이 붙는다. 우선 생각해 볼 군 구조는 다음과 같다.
\begin{equation}
  SU(3)\times SU(2)\times U(1).
\end{equation}
곱셈 기호는 서로 독립적인 인자임을 뜻한다. 하나는 색에, 하나는 약한 이중항에, 나머지 하나는 위상에 작용한다. 이것은 유용한 대략적 그림이지만, 전약 이론 전체를 나타내지는 않는다.\editorialnote{표준모형의 정확한 게이지 군은 관례적으로 \(SU(3)_c\times SU(2)_L\times U(1)_Y\)라고 쓴다. 약한 \(SU(2)_L\)는 왼손잡이 이중항에만 작용하며, 최종적인 전자기 \(U(1)_{\rm em}\)는 대칭 깨짐 뒤 \(SU(2)_L\times U(1)_Y\)에서 나타난다. 강의에서는 처음에 제시한 마지막 두 인자의 동일시가 아직 완전히 일관되지 않는다고 명시한다. 빠져 있는 카이랄 구조와 전약 구조는 뒤 강의들에서 다룬다.}

\subsection{같은 패턴은 렙톤에도 적용된다}

세 색 행 아래에 색이 없는 렙톤 행을 놓자.
\begin{equation}
\begin{array}{c|cc|cc|cc}
 &u&d&c&s&t&b\\ \hline
Q&+\frac23&-\frac13&+\frac23&-\frac13&+\frac23&-\frac13\\
\text{lepton partner}
 &\nu_e&e^-&\nu_\mu&\mu^-&\nu_\tau&\tau^-\\
Q&0&-1&0&-1&0&-1
\end{array}
\end{equation}
서로 이웃한 각 쌍에서는 전하가 한 단위만큼 다르다. 그래서 같은 전하를 띤 약한 보손이 \(u\)와 \(d\), \(c\)와 \(s\), \(t\)와 \(b\), 그리고 전하를 띤 렙톤과 그 중성미자를 연결할 수 있다. \(W\) 꼭짓점은 해당 양자수와 전하 차이에 따라 작용하며, 페르미온이 쿼크인지 렙톤인지는 따지지 않는다.

\begin{figure}[htbp]\centering
\includegraphics[width=0.88\textwidth]{lecture_06_figure_07.png}\caption{완성된 페르미온 표. 위쪽에는 세 색의 쿼크 행, 아래쪽에는 색이 없는 렙톤 행이 있으며 쿼크의 전하와 세 렙톤 이중항이 보인다.}\lectureframe{00:18:55}\end{figure}

뮤온은 전하와 알려진 내재 양자수가 전자와 같지만, 질량은 약 200배 더 크다. 타우는 그보다도 더 무겁다. 중성미자 질량은 매우 작으며, 이 강의에서 논의할 당시에는 자세히 알려져 있지 않았다. 여기까지는 관측된 세대별 질량 패턴을 설명하는 원리가 아직 나오지 않았다.

\begin{classroomqa}\audiencequestion{렙톤, 바리온, 중간자는 왜 그런 이름으로 불리나요?
\lecturetimestamp{00:19:02}}\lecturerresponse{그리스어 어원은 대략 가벼움, 무거움, 중간을 뜻했다. 질량에 근거한 이 명명법은 오해를 낳게 되었다. 뮤온은 질량이 전자와 양성자의 질량 사이에 있어서 한때 뮤 중간자라고 불렸지만, 상호작용을 보면 파이온보다 전자에 훨씬 가까운 렙톤이다. 오늘날 중간자는 쿼크--반쿼크 하드론을 뜻하며 보손이다. 뮤온은 중간자가 아니다.}\end{classroomqa}

\section{약한 상호작용의 수수께끼}

\subsection{구체적인 하전 전류 과정}

중성자 베타 붕괴를 예시 도식으로 살펴보자. 원자가 쿼크의 구성은 다음과 같다.
\begin{equation}
  n=udd,\qquad p=uud.
\end{equation}
다운 쿼크 하나가 다음 전이를 한다.
\begin{equation}
  d\longrightarrow u+W^-,
\end{equation}
그리고 전하를 띤 보손은 다음 입자들을 만든다.
\begin{equation}
  W^-\longrightarrow e^-+\bar\nu_e.
\end{equation}
두 과정을 합치면,
\begin{equation}
  n\longrightarrow p+e^-+\bar\nu_e.
\end{equation}
나머지 원자가 쿼크 두 개는 구경꾼으로서 그대로 이어진다. 첫 번째 꼭짓점에서 중성자는 양성자로 바뀐다. 두 번째 꼭짓점에서는 렙톤 이중항을 연결하는 것과 같은 약한 상호작용으로 \(W^-\)가 관측되는 전자와 전자 반중성미자로 바뀐다.

\subsection{세 결합, 크게 다른 세 가지 인상}

양자전기역학에서 전자 전하 \(e\)는 꼭짓점 진폭이다. 이를 제곱하면 특징적인 방출 확률을 얻으며, 자연 단위계에서는 통상적인 인자들을 모아 다음과 같이 쓴다.
\begin{equation}
  \alpha=\frac{e^2}{4\pi}\simeq\frac1{137}
\end{equation}
강한 상호작용의 결합은 \(g_s\)이며, 전형적인 하드론 에너지 척도에서 이에 대응하는 \(\alpha_s\)는 대략 \(0.2\)일 수 있다. 이는 \(1/137\)보다 훨씬 크므로 쿼크가 충돌할 때 글루온이 흔히 방출된다.

전하를 띤 약한 결합을 \(g_w\)이라 하자. 그 제곱은 터무니없이 작지 않으며, 여기서 사용하는 크기 정도의 추정으로는 전자기 결합과 비슷하다. 그런데도 자유 중성자는 붕괴하기까지 대략 15분을 기다리는 반면, 강한 붕괴는 하드론이 가로지르는 시간 척도에서 일어날 수 있다. 강의의 역사적 추정에는 더 오래된 대략값인 12분이 쓰였지만, 어느 수치를 택해도 같은 수수께끼가 분명히 드러난다.\editorialnote{현재 측정된 자유 중성자의 평균 수명은 약 \(8.8\times10^2\,{\rm s}\)이며, 주요 측정 방법들 사이에는 여전히 실험적 불일치가 있다. 강의에서 제시한 대략적인 수명은 정밀한 값이 아니라 설명에 쓰이는 시간 척도로 유지한다.}\begin{quote}\(g_w\)이 특별히 작지 않다면 약한 과정은 무엇 때문에 약한가?
\end{quote}

답은 꼭짓점에만 있지 않다. 모든 파인만 도표에는 한 꼭짓점에서 다른 꼭짓점까지 중간 전파를 기술하는 인자인 전파인자도 들어 있다.

\section{거리 공간에서 질량을 가진 전파인자}

\subsection{꼭짓점만으로 도표가 완성되지는 않는다}

베타 붕괴 도표에는 약한 꼭짓점 두 개가 있어 \(g_w\)가 두 번 곱해지고, 그 사이를 잇는 \(W\) 선이 있다. 위치 공간 전파인자를 \(\Delta(x-y)\)라 하자. 여기서 \(x\)는 방출 사건의 위치이고 \(y\)는 흡수 또는 변환 사건의 위치다. 이는 두 사건 사이에서 장의 교란이 전달되는 것과 연관된 진폭 커널이다.

정확한 상대론적 전파인자에는 여기서 필요한 것보다 더 많은 구조가 들어 있다. 정성적으로 중요한 특징은 질량에 대한 의존성이다. 짧은 공간꼴 거리에선 거듭제곱형 특이성을 가지며, 개략적으로 다음과 같다.
\begin{equation}
  \Delta(d)\sim\frac1{d^2}.
\end{equation}
질량이 없는 장에는 거듭제곱 꼬리가 남는다. 반면 질량을 가진 장은 분리가 콤프턴 척도보다 커지면 지수적으로 억제된다. 따라서 그래프는 급격히 아래로 꺾인다.

\subsection{자연 단위계로 꺾이는 지점 찾기}

다음과 같이 놓자.
\begin{equation}
  c=\hbar=1.
\end{equation}
그러면
\begin{equation}
  [m]=[E]=[p]=L^{-1}.
\end{equation}
익숙한 관계를 이용하면 각 등식의 의미를 알 수 있다. \(E=mc^2\)는 질량과 에너지를 연결하고, \(p\,x\)의 차원은 \(\hbar\)의 차원이므로 운동량과 거리는 서로 역수의 단위를 갖는다. 따라서 질량 \(m\)으로 만들 수 있는 유일한 길이 척도는 다음과 같다.
\begin{equation}
  d_m\sim\frac1m.
\end{equation}
상수를 복원하면 콤프턴 파장은
\begin{equation}
  d_m\sim\frac{\hbar}{mc},
\end{equation}
가 된다. 교환 입자가 무거울수록 유효 범위는 짧아진다.

\begin{figure}[htbp]\centering
\includegraphics[width=0.78\textwidth]{lecture_06_figure_02.png}\caption{자연 단위계의 차원과 전파인자의 정성적 감소. 검은 곡선은 거듭제곱 꼬리로 이어지고, 빨간 질량 있는 곡선은 역질량 척도 부근에서 아래로 꺾인다.}\lectureframe{00:27:19}\end{figure}

\begin{figure}[htbp]\centering
\begin{tikzpicture}[x=1.0cm,y=0.9cm]
  \draw[->] (0,0) -- (7,0) node[right] {$d$};
  \draw[->] (0,0) -- (0,5) node[above] {$|\Delta(d)|$};
  \draw[thick]
    plot[smooth] coordinates {(0.45,4.5) (0.9,3.3) (1.6,2.5)
      (2.5,1.9) (3.7,1.45) (5.1,1.12) (6.4,0.9)};
  \draw[thick,red]
    plot[smooth] coordinates {(0.45,4.5) (0.9,3.3) (1.6,2.5)
      (2.5,1.9) (3.2,1.55) (3.7,0.95) (4.2,0.36) (4.8,0.1)};
  \draw[dashed] (3.4,0) -- (3.4,1.4);
  \node at (5.7,1.4) {$m=0$};
  \node[red] at (4.75,0.72) {$m>0$};
  \node[below] at (3.4,0) {$d\sim m^{-1}$};
\end{tikzpicture}\caption{칠판의 논증을 보기 쉽게 정리한 도식이다. 4차원 전파인자의 정확한 공식이 아니라 질량이 정하는 범위를 따로 보여준다.}\end{figure}

\begin{classroomqa}\audiencequestion{세로축에는 무엇을 나타냈나요?
\lecturetimestamp{00:28:30}}\lecturerresponse{전파인자의 크기, 즉 \(\Delta\)를 나타낸다. 완전한 진폭, 외부 상태, 위상 공간 인자까지 포함한 뒤 그 제곱이 확률에 기여한다. 이 그래프는 질량 있는 입자를 장거리 교환할 때 진폭이 급격히 사라지는 모습을 보여주려는 것이다.\editorialnote{전파인자만으로는 입자가 고전적 경로를 따라가는 과정에 대해 직접 측정할 수 있는 확률 진폭이 되지 않는다. 강의에서 쓰는 표현은 유용한 범위 추정이다. 전파인자는 전체 산란 또는 붕괴 진폭을 이루는 한 인자다.}}\end{classroomqa}

\section{푸리에 변환과 페르미 한계}

\subsection{실험은 운동량을 측정한다}

실험에서는 내부 선이 시작하고 끝나는 두 시공간점보다 외부 운동량을 측정하는 경우가 보통이다. 베타 붕괴에서 \(W\)가 운반하는 운동량은 중성자 운동량과 양성자 운동량의 차이로 추론하며, 전자와 반중성미자의 전체 운동량과 같다. 위치 공간 전파인자와 운동량 공간 전파인자는 푸리에 변환으로 연결된다.
\begin{equation}
  k=p_n-p_p=p_e+p_{\bar\nu_e}.
\end{equation}

칠판에서 사용한 양의 분모 관례로 쓰면,
\begin{equation}
  \widetilde{\Delta}_W(k)=\frac1{k^2+m_W^2},
  \qquad
  \widetilde{\Delta}_\gamma(k)=\frac1{k^2}.
\end{equation}
질량이 없는 광자 전파인자는 운동량 전달이 작을 때 커진다. 질량이 있는 \(W\) 전파인자는 \(1/m_W^2\) 정도의 척도를 넘지 않는다. 베타 붕괴에서는
\begin{equation}
  |k^2|\ll m_W^2,
\end{equation}
이므로 \(k^2\)의 세부 사항은 무시할 수 있고,
\begin{equation}
  \widetilde{\Delta}_W(k)\simeq\frac1{m_W^2}.
\end{equation}

\begin{classroomqa}\audiencequestion{상대론적 \(k^2\)에는 에너지 항과 공간 운동량 항 사이에 마이너스 부호가 있어야 하나요?
\lecturetimestamp{00:35:11}}\lecturerresponse{그렇다. 공간 성분에 마이너스 부호를 쓰는 민코프스키 계량에서는
\(k^2=(k^0)^2-\boldsymbol{k}^2\)이며, 스칼라 분모는 관례적으로 \(k^2-m_W^2+i\epsilon\)이다. 칠판에서는 이와 동등한 양의 공간꼴 또는 유클리드 크기를 사용한다. 저에너지 베타 붕괴에서는 전달되는 사차원 운동량이 \(m_W\)보다 훨씬 작으므로 어느 표기법을 써도 진폭의 무관한 부호를 제외하면 같은 \(1/m_W^2\) 억제를 준다.}\end{classroomqa}

\subsection{저에너지 진폭}

약한 꼭짓점 두 개와 전파인자 하나를 곱하면
\begin{align}
  \mathcal A_{\rm weak}
  &\sim g_w\,\widetilde{\Delta}_W(k)\,g_w\\
  &\sim\frac{g_w^2}{k^2+m_W^2}
  \simeq\frac{g_w^2}{m_W^2}.
\end{align}
진폭을 제곱하면 특정 붕괴에 맞는 운동학 및 위상 공간 인자를 포함하기 전부터 특징적인 억제가 생긴다.
\begin{equation}
  |\mathcal A_{\rm weak}|^2
  \sim\frac{g_w^4}{m_W^4},
\end{equation}
핵 에너지에서 약한 상호작용이 약한 까닭은 기본 게이지 결합이 터무니없이 작아서가 아니라 매개 입자가 무겁기 때문이다.

\begin{figure}[htbp]\centering
\includegraphics[width=0.86\textwidth]{lecture_06_figure_03.png}\caption{중성자--양성자 약한 상호작용 과정과 내부 선이 운반하는 운동량 전달. 전파인자의 부호를 바로잡는 논의에서 다뤘다.}\lectureframe{00:35:40}\end{figure}

\begin{figure}[htbp]\centering
\includegraphics[width=0.84\textwidth]{lecture_06_figure_04.png}\caption{상자 안의 저에너지 추정식 \(g_w^2/m_W^2\) 옆에 나타낸 베타 붕괴 도식.}\lectureframe{00:37:25}\end{figure}

수치로 보면 이 강의의 정확도에서 \(m_W\)는 대략 양성자 질량의 100배다. 진폭에서 \(m_W^{-2}\)의 인자는 붕괴율에서 \(m_W^{-4}\)의 인자가 된다. 따라서 보통의 게이지 이론에서 볼 수 있는 크기의 결합이라도 수명이 매우 길어질 수 있다.

\begin{classroomqa}\audiencequestion{백만 분의 일 정도로 억제된다는 뜻인가요, 아니면 십억 분의 일 정도인가요? 그리고 \(W\)의 질량은 처음에 어떻게 알려졌나요?
\lecturetimestamp{00:38:11}}\lecturerresponse{강의의 대략적인 계산은 억제되지 않은 강입자 과정에 비해 저에너지 약한 상호작용 확률이 십억 분의 일 규모로 작다는 것을 시사한다. 역사적으로 저에너지 실험에서는 \(g_w\)과 \(m_W\)를 각각 측정하지 않았다. 대신 페르미 상수라고 불리는 작은 조합 \(g_w^2/m_W^2\)을 측정했다. 게이지 결합상수를 전약 이론과 연결하자, 이 조합으로부터 매개자가 매우 무거울 것이라는 예측이 나왔다. 이후 가속기에서 \(W\)를 직접 생성해 질량을 측정했다.}\end{classroomqa}

운동량이 \(m_W\)보다 훨씬 작을 때는 모든 실험에서 거의 같은 상수가 관측된다:
\begin{equation}
  G_F\sim\frac{g_w^2}{m_W^2}.
\end{equation}
이 때문에 페르미는 매개자가 관측되기 훨씬 전인 1930년대에 베타 붕괴를 점입자적인 네 페르미온 상호작용으로 기술할 수 있었다. 충돌 에너지가 높아지면서 실험은 전파자에 감춰진 운동량 의존성에 민감해졌고, 유효 상수가 질량이 큰 매개자 교환의 저에너지 그림자임을 확인했다.\editorialnote{표준 전약 정규화를 적용하면, 트리 수준에서
\[
  \frac{G_F}{\sqrt2}=\frac{g^2}{8m_W^2}.
\]
강의에서는 수치 인자와 스피너 인자를 의도적으로 생략하고 비례 관계 \(G_F\sim g_w^2/m_W^2\)만 남긴다. \(W\)와 \(Z\)는 CERN에서 1983년에 직접 발견되었다. 강의에서 지나가듯 언급한 1970년대 중후반이라는 시기는 직접 관측 시점으로는 너무 이르다.}

약한 상호작용이 무거운 입자에 의해 매개될 수 있다는 생각은 20세기 중반에 나왔다. 강의에서는 이 가능성을 탐구한 이론가 중 줄리언 슈윙거를 언급하고, \(W\)라는 글자가 아마도 weak를 뜻했을 것이라고 덧붙인다. 저에너지 자료만으로도 그런 매개 입자가 핵의 에너지 척도에서 극도로 무거워야 함을 알 수 있었다. 이는 가속기가 입자를 실재 입자로 만들어 낼 만큼 충분한 에너지를 내기 한참 전의 일이었다.

\subsection{붕괴율을 좌우하는 서로 독립적인 세 요소}

핵 결합에 관한 질문을 통해 전이를 좌우하는 논리적으로 구별되는 세 요소를 나눠 볼 수 있다.
\begin{enumerate}
  \item 꼭짓점 결합 상수;
  \item 질량과 운동량 의존성을 포함한 전파인자;
  \item 초기 상태에서 최종 상태로 갈 때 이용할 수 있는 에너지와 위상 공간.
\end{enumerate}
결합 때문에 최종 상태의 에너지가 초기 상태보다 더 커지면 붕괴 채널이 닫힐 수 있다. 이는 운동학적 금지이지, \(g_w\)가 추가로 감소하거나 \(m_W\)가 바뀌는 것이 아니다. 채널이 열려 있어도 이용 가능한 에너지가 적으면 위상 공간에 의해 붕괴가 여전히 크게 억제될 수 있다.

\begin{classroomqa}\audiencequestion{가까이에 있는 양성자의 장이나 결합이 같은 전파인자 메커니즘으로 베타 붕괴를 억제하나요?
\lecturetimestamp{00:41:13}}\lecturerresponse{여기서 논의하는 예에서는 그렇지 않다. 결합이 중요한 이유는 초기와 최종 핵 상태의 에너지를 바꾸기 때문이다. 최종 상태에 도달하는 데 필요한 에너지를 사용할 수 없다면 그 과정은 일어나지 않는다. 결합, 전파인자, 이용 가능한 에너지는 서로 구별되는 세 요소다.}\end{classroomqa}

\begin{classroomqa}\audiencequestion{결합 상수는 정말 상수인가요?
\lecturetimestamp{00:42:55}}\lecturerresponse{고정된 척도에서만 그렇다. 양자 보정에 따라 유효 결합은 운동량이나 거리에 따라 변한다. 여기서는 과정에 관련된 척도에서 각 결합을 다루며, 척도 의존성은 뒤에서 자세히 살펴본다.}\end{classroomqa}

\section{보존 법칙을 버리지 않고 가상 교환 이해하기}

\subsection{무거운 내부 선이 나타날 수 있는 이유}

보통의 중성자 붕괴에는 자유로운 껍질 위 \(W\)를 생성할 만큼 에너지가 없다. 도표에 그려진 \(W\)는 내부 선이며 껍질에서 벗어나 있다. 강의에서는 이를 전통적인 에너지--시간 관계로 소개한다.
\begin{equation}
  \Delta E\,\Delta t\lesssim\hbar:
\end{equation}
고전적으로 에너지가 맞지 않는 중간 상태는 \(\Delta E\)만큼의 차이가 있을 때 대략 \(\hbar/\Delta E\) 정도의 시간 동안만 기여할 수 있다.

터널링 비유를 들면 의도한 그림을 쉽게 떠올릴 수 있다. 원자핵의 퍼텐셜 우물에 알파 입자가 있다고 하자. 그 에너지는 장벽 꼭대기보다 낮아서 고전적으로는 빠져나올 수 없다. 양자역학적으로 파동함수는 장벽 너머까지 뻗으며 입자는 처음과 같은 전체 에너지로 반대편의 한 지점에 나타날 수 있다. 장벽 내부를 보통의 고전적 궤적으로 설명하려 하면 퍼텐셜 에너지가 너무 큰 것처럼 보이고 운동 에너지는 음수여야 한다. 이렇게 실패하는 고전적 그림은 내부 \(W\)를 실재 입자로 보려는 것과 비슷하다.

\begin{figure}[htbp]\centering
\includegraphics[width=0.78\textwidth]{lecture_06_figure_06.png}\caption{터널링 비유. 시작점과 반대쪽 경사면에 표시된 점의 전체 에너지는 같으며, 그사이 장벽은 고전적으로 접근할 수 없다.}\lectureframe{00:46:00}\end{figure}

\begin{figure}[htbp]\centering
\begin{tikzpicture}[x=1.0cm,y=0.9cm]
  \draw[->] (0,0) -- (7,0) node[right] {$x$};
  \draw[->] (0,0) -- (0,4.8) node[above] {$V(x)$};
  \draw[thick]
    plot[smooth] coordinates {(0.3,1.4) (1.0,0.7) (1.7,0.55)
      (2.1,0.9) (3.2,3.9) (4.2,2.1) (5.0,0.75) (6.6,0.18)};
  \draw[dashed] (1.7,0.75) -- (5.0,0.75);
  \fill (1.7,0.75) circle (0.06);
  \fill (5.0,0.75) circle (0.06);
  \node at (3.35,0.35) {$\Delta E\,\Delta t\lesssim\hbar$};
\end{tikzpicture}\caption{장벽과 같은 에너지의 두 지점을 알아보기 쉽게 재구성했다.}\end{figure}

베타 붕괴 그림에서 중성자는 양성자와 내부 \(W^-\)로 요동할 수 있지만, 무거운 선은 곧 사라져야 한다. 보손은 다시 흡수되어 중성자가 되돌아올 수도 있고, 두 번째 꼭짓점에서 끝나면서 훨씬 가벼운 전자와 반중성미자를 만들 수도 있다. 초기와 최종 상태의 에너지는 같다.\editorialnote{파인만 꼭짓점마다 에너지와 운동량은 보존된다. 가상 입자가 에너지를 문자 그대로 빌렸다가 나중에 갚는 것은 아니다. 정확히는 내부 선의 사차원 운동량에 대해 적분하는데, 그 값은 실재 입자의 질량 껍질 관계 \(k^2=m_W^2\)를 만족하지 않아도 된다. 불확정성 원리와 터널링이라는 표현은 강의의 설명 방식을 그대로 살리기 위해 유지했지만, 현대 양자장론에서는 껍질에서 벗어난다는 것이 정확한 정식화다.}

\subsection{실재 W 보손}

충분히 에너지가 높은 가속기에서는 실재 \(W\)를 만들 수 있다. 이 입자는 가상 내부 선의 가정된 불확정성 시간과는 무관한 고유 수명을 갖는다. 부스트된 실재 \(W\)는 보통의 상대론적 시간 지연에 따라 실험실계에서 더 오래 산다. 정지한 \(W\)가 렙톤 두 개로 붕괴하는 2체 과정에서는
\begin{equation}
  W^-\longrightarrow e^-+\bar\nu_e,
\end{equation}
전자와 반중성미자가 큰 운동 에너지를 가지며 거의 서로 반대 방향으로 나아간다. 각각의 에너지는 대략 \(m_W/2\)이다. 정지 에너지가 사라진 것이 아니라 최종 상태의 에너지로 바뀐 것이다. 뮤온과 타우 채널뿐 아니라 쿼크 채널도 가능하다.

\begin{classroomqa}\audiencequestion{실재 \(W\)는 얼마나 오래 사나요? 그 수명을 같은 결합으로 계산할 수 있나요?
\lecturetimestamp{00:49:49}}\lecturerresponse{그렇다. 허용된 각 붕괴의 진폭을 계산하고, 이를 제곱한 뒤 최종 상태 위상 공간에 대해 적분하고 부분 붕괴폭을 합한다. 전체 붕괴폭의 역수가 수명이다. 녹취에서 말한 수치 추정은 알아듣기 어렵지만, 고유 수명이 극도로 짧고 움직이는 \(W\)에는 통상적인 시간 지연이 적용된다는 정성적 요점은 분명하다.\editorialnote{측정된 붕괴폭 \(\Gamma_W\simeq2.1\,{\rm GeV}\)를 사용하면 수명은 \(\tau_W=\hbar/\Gamma_W\simeq3.2\times10^{-25}\,{\rm s}\)이다. 이 현대적 수치는 강의 내용으로 돌리지 않으면서 손상된 수치 대화를 바로잡아 준다.}}\end{classroomqa}

\begin{classroomqa}\audiencequestion{진폭을 제곱하면 어떻게 붕괴율이 되고, 이어서 수명을 구하나요?
\lecturetimestamp{00:50:52}}\lecturerresponse{행렬 원소의 제곱에 이용 가능한 최종 상태의 밀도, 보존 델타 함수, 정규화 인자를 결합하면 단위 시간당 확률인 붕괴폭 \(\Gamma\)을 얻는다. 생존 확률은 \(e^{-\Gamma t}\)이다. \(\hbar=1\)인 단위계에서는 \(\tau=1/\Gamma\)이고, 단위를 복원하면 \(\tau=\hbar/\Gamma\)이다. 따라서 약한 붕괴폭이 작다는 것은 약한 상호작용의 수명이 길다는 뜻이다.}\end{classroomqa}

\begin{classroomqa}\audiencequestion{장벽을 터널링하는 입자에 속도를 부여할 수 있나요?
\lecturetimestamp{00:52:42}}\lecturerresponse{추론된 거리를 추론된 이동 시간으로 나눌 수는 있지만, 그렇게 얻은 양을 고전적 입자나 신호의 속도로 해석해서는 안 된다. 장벽 내부 궤적을 측정하려 하면 상태가 바뀌고 원래의 터널링 실험은 파괴된다.}\end{classroomqa}

\section{가상 입자를 관측한다는 것은 무슨 뜻인가?}

\subsection{외부 입자와 내부 선}

실재 입자는 상호작용 영역에서 멀리 떨어진 곳에서 준비하거나 검출할 수 있는 입사 및 출사 점근 상태다. 가상 입자는 진폭의 섭동론적 표현에 나타나는 내부 선이다. 검출기에 추가로 숨겨진 흔적이 생기는 것은 아니다.

\begin{classroomqa}\audiencequestion{실재 입자와 가상 입자의 차이를 명확하게 정의할 수 있나요?
\lecturetimestamp{00:55:05}}\lecturerresponse{실재 입자는 상호작용 영역으로 들어오거나 그 영역을 빠져나가는 외부 껍질 위 상태다. 가상 입자는 파인만 도표의 내부 껍질 밖 선이다. 후자는 수학적 전개의 일부이지, 독립적으로 관측되는 대상이 아니다.}\end{classroomqa}

\(\Delta t\)만큼만 지속되는 과정을 분해해서 보려면 탐침에 \(1/\Delta t\) 정도의 주파수가 필요하며, 따라서 에너지는 \(\hbar/\Delta t\) 정도여야 한다. 높은 장벽 아래에 있는 알파 입자를 분해할 만큼 예리한 광자는 입자를 장벽 영역까지 끌어올릴 에너지를 공급할 수 있다. 마찬가지로 수명이 짧은 내부 \(W\)를 분해할 만큼 에너지가 큰 탐침은 실재 \(W\)를 만들 수 있다. 관측을 시도하면 원래 과정은 다른 실험으로 바뀐다.

강의에서는 이것을 양자역학에서 되풀이되는 난처한 상황으로 표현한다. 금지된 중간 상태 그림이 맞는지 확인하는 데 필요한 개입이, 겉보기에 금지된 그 상태를 실재하게 만드는 바로 그 에너지를 공급한다.

\subsection{양성자의 요동}

섭동론에서 양성자 상태에는 다음과 같은 구성의 진폭이 포함된다.
\begin{equation}
  p\longleftrightarrow n+W^+,
\end{equation}
중성자에도 마찬가지로 다음이 있다.
\begin{equation}
  n\longleftrightarrow p+W^-.
\end{equation}
다른 상호작용이 개입하지 않으면 내부 보손은 다시 흡수된다. 강한 탐침은 대신 \(W\)를 산란 후 입자들 가운데 남길 수 있다. 이를 요동을 걷어내거나 \(W\)를 만드는 행위라고 부를 수 있다. 표현은 달라도 측정 가능한 진술, 즉 같은 최종 상태의 진폭을 말한다.

\begin{classroomqa}\audiencequestion{강한 탐침이 산란 후 \(W\)를 만들면, 이미 있던 입자를 걷어낸 건가요, 아니면 새 입자를 만든 건가요?
\lecturetimestamp{01:01:20}}\lecturerresponse{두 그림 모두 수학을 정리하는 데 쓸 수 있다. 한 시간 순서 기술에서는 양성자가 \(n+W^+\)로 요동하고 보손이 다시 흡수되기 전에 탐침이 중성자를 떼어낸다. 다른 기술에서는 충돌이 최종 \(W\)를 만든다. 불변인 것은 관측된 최종 상태의 계산된 확률이지, 측정 전에 무엇이 존재했는지에 대한 유일한 고전적 이야기 방식이 아니다.}\end{classroomqa}

충돌 에너지를 높여도 표적 내부에 숨겨져 있던 물체의 고정된 목록이 드러나는 것은 아니다. 대신 보존 양자수가 허용하는 가능한 최종 상태가 많이 열린다. \(W\) 하나, 여러 보손, 쿼크 제트, 하드론, 그 밖의 입자들이 여기에 포함된다. 기본 장이라도 충분히 세게 충돌시키면 양성자와 다른 복합 물체를 만들어낼 수 있다. 산란 후 나오는 입자가 충돌 전에 구성 성분으로 존재했을 필요는 없다. 강의에서는 이 점을 우스꽝스러울 정도로 극단적인 예까지 밀고 나간다. 에너지와 양자수가 허용한다면 확률이 엄청나게 낮더라도 탁자나 의자, 수소 원자조차 논리적으로 배제되지 않는다. 이 농담에는 충돌 잔해를 단순히 세어 기본성을 정의하지 말라는 진지한 경고가 담겨 있다.

\subsection{진공 요동과 온도}

양성자에 붙어 있는 요동은 진공 도표와 같지 않다. 전자는 상호작용하는 양성자 상태의 성분이고, 진공 요동에는 외부 양성자 선이 없다. 어느 경우든 충분히 에너지가 높은 탐침은 허용된 최종 상태에 해당하는 실재 입자를 만들 수 있다. 예를 들어 질량 중심 에너지가 문턱값을 넘으면 에너지가 높은 광자 두 개가 충돌해 전자--양전자 쌍을 만들 수 있다.

\begin{classroomqa}
\audiencequestion{이 반응들은 온도에 의존하나요?
\lecturetimestamp{01:05:17}}
\lecturerresponse{기본 진폭을 만드는 데 열적 열원이 필요한 것은 아닙니다.
온도는 통계적으로 중요합니다. 온도가 높은 계에는 운동 에너지가 큰 입자가 많으므로, 문턱을 넘는 충돌이 흔해집니다. 결합된 알파 입자가 있는 계를 가열하면 양자 터널링에 더해 장벽을 넘어 탈출하기에 충분한 에너지가 공급될 수도 있습니다. 고립된 고에너지 입자 몇 개가 이루는 것은 빔이고, 열적으로 분포한 다수의 고에너지 입자가 이루는 것은 뜨거운 열원입니다.}
\end{classroomqa}

\begin{classroomqa}
\audiencequestion{양성자 요동에서 방출되는 \(W\)는 진공의 요동인가요?
\lecturetimestamp{01:06:56}}
\lecturerresponse{아닙니다. \(p\leftrightarrow n+W^+\)에서 내부 \(W\)는 상호작용하는 양성자 상태의 요동입니다.
외부 양성자 없이 닫힌 가상 과정이 일어나면 그것은 진공 요동입니다. 두 경우 모두 같은 장론 규칙을 따르지만 서로 다른 진폭에 해당합니다.}
\end{classroomqa}

\begin{classroomqa}
\audiencequestion{렙톤 \(W\) 꼭짓점에 반중성미자가 필요한 이유는 무엇인가요?
\lecturetimestamp{01:08:51}}
\lecturerresponse{반중성미자는 보존되는 렙톤 양자수를 제공하며, \(e\), \(\nu_e\), \(W\)를 연결하는 기본 하전 전류 꼭짓점의 일부입니다.
위에 막대가 붙은 표기는 방출된 전자 반중성미자와 중성미자를 구별합니다.}
\end{classroomqa}

\section{에너지 변환은 어디까지 가능할까?}

다음으로 강의에서는 에너지가 입자를 만들 수 있다는 규칙을 의도적으로 터무니없는 한계까지 밀어붙인다. 커피콩 하나가 초콜릿 케이크가 될 수 있을까? 그런 현실적인 과정은 없지만, 상대론적 양자 이론은 최종 물질이 처음 물체와 닮아야 한다는 규칙을 두지 않는다. 충분한 에너지와 올바른 보존 전하, 그리고 영이 아닌 진폭이 필요하다. 질량이 킬로그램 정도인 물체의 에너지는 대략 \begin{equation}
  Mc^2\sim9\times10^{16}\ {\rm J},
\end{equation}이므로, 특정한 킬로그램 규모의 물체를 만들어 내는 미시적 충돌은 상상하기 어려울 만큼 일어날 가능성이 낮다.

플랑크 질량,
\begin{equation}
  m_{\rm P}=\sqrt{\frac{\hbar c}{G}}
  \simeq2.2\times10^{-5}\ {\rm g},
\end{equation}
은 가능한 질량의 최댓값이 아닙니다. 양자 입자가 콤프턴 파장으로 국소화될 때 중력이 강해지는 규모입니다. 킬로그램 하나에는 플랑크 질량이 대략 \(5\times10^7\)개 들어 있습니다.

\begin{classroomqa}
\audiencequestion{충돌 에너지가 플랑크 규모를 넘으면 의자나 행성처럼 질량이 임의로 큰 물체가 생길 수 있나요?
\lecturetimestamp{01:10:37}}
\lecturerresponse{작은 영역에 집중된 질량 중심 에너지가 플랑크 규모를 넘으면, 압도적으로 자연스러운 결과는 중력 붕괴, 즉 블랙홀입니다. 고도로 조직된 거시적 물체가 생기는 것은 극히 드문 요동일 것입니다. 강의에서는 시간 반전을 떠올립니다. 의자가 블랙홀에 빠질 수 있으므로 그 역방향 진폭은 정확히 0은 아닙니다. 그 확률은 엄청난 엔트로피 인자만큼 억제되므로, 가능성과 실제로 일어날 법한 정도를 혼동해서는 안 됩니다.}
\end{classroomqa}

이 이야기는 앞서 한 경고를 다시 강조합니다. 양자장론은 지정된 초기 상태와 최종 상태 사이의 진폭을 예측합니다. 충돌 전부터 모든 생성물이 표적 안에 문자 그대로 들어 있었다는 식의 그림 같은 주장은 이 예측에 포함되지 않습니다.

\section{보손과 그 상호작용 및 통계}

\subsection{광자는 산란할 수 있다}

보손이라고 해서 상호작용하지 않는 것은 아니다. 보손인지 페르미온인지는 교환 통계를 규정하며, 상호작용 꼭짓점은 별개의 동역학적 사실이다. 양자전기역학에는 나무 수준의 네 광자 꼭짓점이 없고 저에너지 광자--광자 산란도 극히 작기 때문에, 고전적인 빛살은 보통 서로를 통과한다. 그럼에도 두 광자는 가상 전자--양전자 고리를 통해 상호작용할 수 있다:
\begin{equation}
  \gamma+\gamma
  \longrightarrow(e^-e^+)_{\rm virtual}
  \longrightarrow\gamma+\gamma.
\end{equation}
질량중심 에너지가 충분히 높으면, 대신 실제 입자 쌍을 만들 수도 있다.
\begin{equation}
  \gamma+\gamma\longrightarrow e^-+e^+.
\end{equation}
광학 광자의 에너지는 수 전자볼트인 반면, 쌍 생성 문턱 에너지는 대략 \(2m_e c^2\approx1.02\,{\rm MeV}\)이다. 따라서 보통의 광학 에너지에서는 가상 과정을 통한 상호작용이 크게 억제되며, 이 때문에 교차하는 빛살은 서로를 알아차리지 못하는 것처럼 보인다.

\begin{classroomqa}
\audiencequestion{가상 \(W\) 보손은 보손이라는 이유만으로 다른 보손과 상호작용하나요?
\lecturetimestamp{01:12:47}}
\lecturerresponse{상호작용할 수는 있지만, 보스 통계를 따른다는 이유만으로 그러는 것은 아닙니다. 동역학적 상호작용은 게이지 꼭짓점으로 정해집니다. 통계는 동일한 입자들을 교환할 때 다입자 진폭이 어떻게 변하는지를 정합니다. 광자는 전하를 띤 가상 입자를 통해 약하게 상호작용하며, 전약 이론에는 비아벨 \(W\) 보손의 직접적인 자기상호작용도 있습니다.}
\end{classroomqa}

\subsection{보스 증강은 힘이 아니다}

여러 개의 동일한 입자가 서로 가까운 양자 상태에 있을 때 페르미온과 보손의 통계 차이가 두드러집니다. 보손은 한 상태에 여러 입자가 함께 점유하는 경향이 있고, 페르미온은 동일한 입자들이 같은 상태를 점유하지 못하게 합니다. 전형적인 고에너지 충돌에서는 입자들이 서로 충분히 떨어진 운동량 상태로 날아가므로, 레이저나 응축체에서만큼 이 통계적 효과가 뚜렷하지 않을 수 있습니다.

원칙적으로 \(W\) 보손에도 같은 보스 증강이 있습니다. 강의에서는 \(W\) 보손으로 레이저를 만드는 계를 상상해 볼 수 있다고 농담합니다. 보손의 엄청난 질량과 극도로 짧은 수명 때문에 그런 장치는 상상할 수 없을 만큼 비현실적이지만, 보손의 통계가 보스 통계인 것은 맞습니다.

\begin{classroomqa}
\audiencequestion{양성자나 중성자는 무슨 색인가요?
\lecturetimestamp{01:18:41}}
\lecturerresponse{물리적 강입자는 색 단일항이며, 흔히 무색 또는 흰색이라고 부릅니다. 간단히 말하면 중입자는 빨강, 초록, 파랑 쿼크를 하나씩 포함합니다. 더 정확히 말해 색 파동함수는 다음과 같은 반대칭 중첩입니다.
\[
  |{\rm singlet}\rangle
  =\frac1{\sqrt6}\epsilon_{abc}|q^a q^b q^c\rangle,
\]
이는 이름을 붙인 각 쿼크에 고정된 색을 고전적으로 지정하는 것이 아닙니다.}
\end{classroomqa}

\begin{classroomqa}\audiencequestion{이 강의에서 보스 응축을 다루나요?\lecturetimestamp{01:19:25}}\lecturerresponse{네, 힉스 현상을 통해 다룹니다. 진공 전체에 걸쳐 영이 아닌 힉스장이 존재한다는 점은 중요한 측면에서 보손 응축과 유사합니다. 자발적 대칭 깨짐을 소개한 뒤에 이 비유를 살펴보겠습니다.\editorialnote{힉스 진공 기대값은 한 상태를 점유하는 질량껍질 위 힉스 입자들로 이루어진 보통의 기체를 문자 그대로 뜻하지는 않습니다. 응축은 결맞은 보손장과 선택된 진공을 이해하는 데 유용한 비유이지, 미시적 실체를 완전히 동일시하는 설명은 아닙니다.}}\end{classroomqa}

\section{빠져 있는 중성 약력 보손}

이제 군 표로 돌아가 봅시다. 수평 방향의 약한 대칭이 \(SU(2)\)라면,
\begin{equation}
  \dim SU(2)=3
\end{equation}
개의 생성원이 있으므로 게이지장이 세 개 있습니다. 하전 조합인 \(W^+\)와 \(W^-\)는 이 가운데 두 개만 설명합니다. 그러므로 세 번째 중성장이 필요합니다.
이 장은 관측된 \(Z\) 보손과 밀접한 관련이 있지만, 단순히 광자인 것은 아닙니다. 중성 조합을 정확히 표현하려면 세 번째 \(SU(2)_L\) 장과 \(U(1)_Y\) 장을 함께 고려해야 합니다. 이들 사이의 혼합과 전하의 올바른 정의는 일부러 다음 강의로 미뤄 둡니다.

\section{자발적 대칭 깨짐}

\subsection{대칭적인 선택에도 불안정한 망설임이 있을 수 있다}

원형 저녁 식탁을 떠올려 봅시다. 각 자리에는 왼쪽과 오른쪽에 식기가 하나씩 놓여 있습니다. 아무도 움직이기 전에는 왼쪽과 오른쪽이 동등합니다. 한 손님이 오른쪽으로 손을 뻗으면, 이웃 사람들의 선택도 모두 정해져 식탁 전체가 오른손잡이 배치로 자리 잡습니다. 규칙은 대칭이었지만, 실제로 구현된 배열은 두 가능성 중 하나를 골랐습니다. 이 이야기는 그저 비유일 뿐이지만, 우리가 필요한 요소들을 잘 드러냅니다. 대칭성, 서로 동등한 여러 결과, 그리고 아주 작은 교란이 전체적으로 한 결과를 선택하게 하는 불안정성입니다.

\subsection{앞면--뒷면 격자}

더 단순한 모형을 위해 격자의 각 자리에 두 상태 변수 \(s_i=\pm1\)를 둡시다. \(+1\)은 앞면, \(-1\)는 뒷면이라고 합시다. 이웃한 자리들이 같은 방향으로 정렬되기를 선호한다고 합시다.
\begin{equation}
  H=-J\sum_{\langle ij\rangle}s_i s_j,
  \qquad J>0.
\end{equation}
평행한 이웃은 반평행한 이웃보다 에너지가 낮습니다. 따라서 바닥 상태가 두 개 퇴 degeneracy를 이룹니다.
\begin{equation}
\begin{aligned}
  s_i&=+1\quad\text{for every }i,
  \\
  s_i&=-1\quad\text{for every }i.
\end{aligned}
\end{equation}
해밀토니안은 전체 반전 \(s_i\to-s_i\)에 대해 불변이지만, 둘 중 어느 한 바닥 상태를 선택하면 그 상태는 불변이 아닙니다.

이제 아주 멀리 떨어진 경계의 스핀 하나를 앞면으로 고정합시다. 보통 단거리 상호작용 계에서는 멀리 떨어진 자유도 하나가 이곳에 영향을 주지 않으리라 예상할 수 있습니다. 하지만 질서화된 상에서는 그 경계 조건이 앞면만 있는 진공을 선택합니다. 일단 그 진공이 선택되면, 앞면이 가득한 영역에 뒷면 하나가 생길 때 이웃과 방향이 달라 에너지가 듭니다. 국소 들뜸은 바탕 방정식에 앞면--뒷면 대칭이 남아 있더라도 선택된 바닥 상태를 기준으로 정의됩니다.

\subsection{명시적 대칭 깨짐과 영역벽}

명시적으로 대칭을 깨려면 국소적인 편향 항을 더하면 됩니다.
\begin{equation}
  H_{\rm explicit}
  =-J\sum_{\langle ij\rangle}s_i s_j-h\sum_i s_i,
  \qquad h\ne0.
\end{equation}
그러면 한 자리만 놓고 보아도 앞면과 뒷면의 에너지가 같지 않아, 선호되는 진공은 하나뿐입니다. 자발적 대칭 깨짐에서는 \(h=0\)입니다. 국소 규칙은 대칭을 유지하고 두 진공은 여전히 축퇴되어 있지만, 계는 둘 중 하나에 자리 잡습니다.

대신 가장 왼쪽 경계는 모두 앞면, 가장 오른쪽 경계는 모두 뒷면이라고 합시다. 최저 에너지 배치는 두 진공 사이를 이어야 합니다. 이 경계면이 영역벽입니다.

\begin{figure}[htbp]\centering
\includegraphics[width=0.9\textwidth]{lecture_06_figure_08.png}\caption{앞면--뒷면 격자를 점선 영역벽이 가로지른다. 서로 반대인 경계 조건 때문에 계에는 두 축퇴된 질서 상태 사이의 경계면이 생긴다.}
\lectureframe{01:30:35}
\end{figure}

\begin{classroomqa}
\audiencequestion{자발적 대칭 깨짐과 보통의 명시적 대칭 깨짐은 어떻게 다른가요?
\lecturetimestamp{01:28:36}}
\lecturerresponse{명시적 대칭 깨짐에서는 국소 해밀토니안이 이미 앞면과 뒷면에 서로 다른 에너지를 부여하므로, 대칭으로 연결된 바닥 상태 두 개가 없습니다. 자발적 대칭 깨짐에서는 방정식이 대칭을 유지하고 진공들은 축퇴되어 있으며, 경계 조건이나 임의로 작은 교란이 그중 하나를 선택합니다. 서로 다른 진공 사이를 잇는 영역벽이 가능하다는 점은 후자의 상황을 판별하는 기준입니다.}
\end{classroomqa}

이 모형은 의도적으로 매우 단순하지만, 힉스장에 적용할 핵심 논리를 담고 있습니다. 대칭적인 이론, 불안정하거나 축퇴된 진공 구조, 선택된 바닥 상태, 그리고 그 선택에 따라 정해지는 물리적 결과입니다.

\section{이 강의에서 배운 내용}

페르미온 표는 표준모형의 군을 물리적 방향들의 집합으로 바꾸어 보여 줍니다. 색 \(SU(3)\)은 빨강, 초록, 파랑 사이를 수직으로 이동시키고, 하전 약력 변환은 쿼크와 렙톤 이중항 사이를 수평으로 이동시킵니다. 위상 대칭은 예비적인 \(U(1)\) 인자를 제공합니다. 게이지 보손 방출은 이러한 군 변환을 상호작용으로 구현합니다.

따라서 베타 붕괴의 약함은 특별히 작은 \(g_w\) 때문이 아니라 무거운 매개자 때문이라고 설명할 수 있습니다. 거리 공간에서 질량은 콤프턴 척도 \(\hbar/(m_Wc)\) 너머로 전파자를 차단합니다. 운동량 공간에서는 핵자 운동량에서
\begin{equation}
  \mathcal A_{\rm weak}
  \sim\frac{g_w^2}{k^2+m_W^2}
  \simeq\frac{g_w^2}{m_W^2}
\end{equation}
이 됩니다. 과거의 저에너지 실험은 이 조합을 페르미 상수로 측정했습니다.

수업 중 이어진 논의는 이 공식이 뜻하는 것과 뜻하지 않는 것을 분명히 합니다. 내부 \(W\) 선은 가상이며 질량껍질에서 벗어나 있지만, 에너지--운동량은 여전히 보존됩니다. 충분히 강한 탐침은 실재 \(W\)를 만들 수 있지만, 그렇게 되면 다른 최종 상태 과정이 생깁니다. 충돌 생성물은 미리 존재하던 구성 성분의 목록이 아닙니다. 온도는 고에너지 입사 입자의 수를 바꾸며, 광자는 전하를 띤 가상 입자를 통해 상호작용할 수 있고, 보스 통계는 동역학적 힘과 구별됩니다.

마지막으로 \(SU(2)\)에는 빠져 있는 중성 게이지장이 필요한데, 이 장은 \(Z\)와 관련이 있습니다. 앞면--뒷면 격자는 자발적 대칭 깨짐을 보여 주는 첫 모형입니다. 다음 단계에서는 아직 마무리되지 않은 이 두 주제, 즉 전약 혼합과 힉스 메커니즘을 결합합니다.
